\documentclass[11pt,a4paper]{article}

% --- Packages ---
\usepackage[utf8]{inputenc}
\usepackage[T1]{fontenc}
\usepackage[english]{babel}
\usepackage{geometry}
\geometry{a4paper, margin=2.5cm}
\usepackage{xcolor}
\usepackage{titlesec}
\usepackage{titling}
\usepackage{fancyhdr}
\usepackage{hyperref}
\usepackage{enumitem}
\usepackage{graphicx}
\usepackage{listings}
\usepackage{tcolorbox}

% --- Colors ---
\definecolor{navy}{HTML}{002060}
\definecolor{silver}{HTML}{C0C0C0}
\definecolor{lightgray}{HTML}{F5F5F5}

% --- Header/Footer ---
\pagestyle{fancy}
\fancyhf{}
\rhead{\small\textcolor{gray}{Euro-Shield 6Y Horizon - Project Methodology}}
\cfoot{\small\textcolor{gray}{\thepage}}
\renewcommand{\headrulewidth}{0.4pt}
\renewcommand{\footrulewidth}{0pt}

% --- Title Styling ---
\titleformat{\section}{\color{navy}\normalfont\Large\bfseries}{\thesection}{1em}{}
\titleformat{\subsection}{\color{navy}\normalfont\large\bfseries}{\thesubsection}{1em}{}

% --- Metadata ---
\title{
    \vspace{2cm}
    \textbf{\Huge\color{navy}Project Methodology Documentation} \\
    \vspace{1cm}
    {\Large Euro-Shield 6Y Horizon --- Structured Product Pricing} \\
    \vspace{0.5cm}
    {\large SKEMA Business School --- Financial Markets \& Investments}
    \vspace{2cm}
}
\author{\textbf{Quentin LY}}
\date{April 2026}

% --- Document Start ---
\begin{document}

\maketitle
\thispagestyle{empty}
\vspace{2cm}
\begin{center}
    \rule{\linewidth}{0.5pt}
\end{center}
\vfill
\newpage

\tableofcontents
\newpage

\section{Overview}
The objective of this project is to build an institutional-grade pricing and analysis engine for the \textit{Euro-Shield 6Y Horizon} structured product. This document outlines the technical environment, the project architecture, and the modular methodology used to decompose this complex financial instrument into verifiable ``Lego'' blocks.

\section{Technical Environment}
The project is built using a modern Python ecosystem focused on performance, reproducibility, and clarity.

\subsection{Core Tools}
\begin{itemize}
    \item \textbf{Package \& Env Manager:} \texttt{uv} (v0.11+). Used for ultra-fast dependency resolution and virtual environment isolation.
    \item \textbf{Language:} Python 3.9+.
    \item \textbf{Documentation/Reporting:} LaTeX and Typst. Used for generating high-quality PDF reports and presentations with minimal overhead.
\end{itemize}

\subsection{Primary Dependencies}
\begin{itemize}
    \item \texttt{pandas}: Data manipulation and results tabulating.
    \item \texttt{matplotlib}: Institutional-grade technical plots and payoff diagrams.
    \item \texttt{scipy}: Financial optimization and normal distribution calculations.
    \item \texttt{pymupdf}: PDF parsing for raw document analysis.
    \item \texttt{tabulate}: Formatting console outputs for README generation.
\end{itemize}

\section{Project Structure}
The repository is organized following a strict data-science pipeline structure to ensure separation of concerns between raw inputs, processed data, and visual assets.

\subsection{Directory Hierarchy}
\begin{itemize}[label=--]
    \item \texttt{data/raw/}: Contains original case study documents (PDFs) and market data (Excel files).
    \item \texttt{data/processed/}: Contains intermediate calculation results (CSV, MD) used to populate the main README.
    \item \texttt{images/}: Stores all generated plots, diagrams, and ``Lego'' decompositions.
    \item \texttt{src/}: The core logic directory containing modular scripts for each step of the pricing workflow.
    \item \texttt{methodology/}: Contains documentation sources (Typst/LaTeX) and final PDF reports.
\end{itemize}

\subsection{Key Configuration Files}
\begin{itemize}[label=--]
    \item \texttt{pyproject.toml}: Defines the project workspace, dependencies, and environment requirements.
    \item \texttt{README.md}: The central hub of the project, summarizing all findings.
    \item \texttt{.python-version}: Locks the local environment to a specific Python version.
\end{itemize}

\section{Methodology: Step-by-Step Build}
The project follows a 9-step methodology, moving from theoretical payoff formalization to secondary market dynamics.

\subsection{Step 1: Payoff Formalization \& Visualization}
Decomposition of the product into ``Lego'' blocks: a Long Zero-Coupon Bond, a Long ATM Call, and a Short OTM Call (Cap). Mathematical verification that the sum of these blocks equals the promised product payoff.

\subsection{Step 2: Yield Curve Bootstrapping}
Conversion of raw market swap rates into a continuous zero-coupon yield curve. This enables precise discounting for the capital guarantee component.

\subsection{Step 3 \& 4: Black-Scholes Pricing \& Margin}
Implementation of the Black-Scholes model to price the embedded options. The difference between the selling price (€100) and the manufacturing cost determines the bank's upfront margin.

\subsection{Step 5: Greeks \& Risk Exposure}
Calculation of Delta, Gamma, Vega, Rho, and Theta. This step identifies the product as a ``Short Vega'' and ``Long Theta'' instrument.

\subsection{Step 6: Sensitivity Analysis}
Stress-testing the bank margin against market shifts (Rates $\pm$ 50bp, Volatility $\pm$ 1\%, Maturity shifts). This quantifies the robustness of the pricing model.

\subsection{Step 7: Secondary Market Pricing}
Valuation of the product after inception. Analysis of the ``Pull-to-par'' effect and the impact of funding penalties on bid prices during client buybacks.

\section{Conclusion}
By maintaining a modular codebase in \texttt{src/} and a central documentation hub in \texttt{README.md}, the project ensures that every financial assumption is transparent, verifiable, and reproducible. The use of \texttt{uv} ensures that the technical environment is easy to reconstruct on any machine.

\end{document}
